\section{Product Interface}
\label{sec:interface}

\cs\ is organised as a persistent left-hand navigation over twelve
screens that share one visual system---a single accent colour, a common
metric-tile and panel component, and an always-visible status strip that
shows the Docker and local-model state. This section walks through the
screens a developer meets during a full audit; all figures show the
platform analysing the case-study project of \cref{sec:results}.

\subsection{Initialisation}
On launch, \cs\ runs a four-stage readiness check
(\cref{fig:boot}): host resources, Docker runtime, provisioning of the
local model container, and the model pull itself. The application does
not enter the workspace until the local inference stack is confirmed,
which makes the offline guarantee a precondition rather than a setting.

\begin{figure}[t]
  \centering
  \begin{subfigure}{0.485\columnwidth}
    \includegraphics[width=\linewidth]{boot-1}
    \caption{Hardware verification}
  \end{subfigure}\hfill
  \begin{subfigure}{0.485\columnwidth}
    \includegraphics[width=\linewidth]{boot-3}
    \caption{Model container pulse}
  \end{subfigure}
  \caption{Initialisation sequence. \cs\ confirms the local Docker and
           \mbox{Llama~3.2} stack before opening the workspace.}
  \label{fig:boot}
\end{figure}

\subsection{Dashboard}
The dashboard (\cref{fig:dashboard}) is the audit summary: key-performance
tiles (file count, total findings, mean complexity, build state), a
severity-distribution ring, a panel of model-generated insights, and a
short natural-language security summary. It answers ``how is this
project doing'' in one view before the developer drills in.

\begin{figure}[t]
  \centering
  \includegraphics[width=\columnwidth]{dashboard}
  \caption{Dashboard---audit summary with severity distribution and
           model-generated insight cards.}
  \label{fig:dashboard}
\end{figure}

\subsection{Repository and Configuration}
The repository screen registers a project from a Git URL or a single
source file and lists the workspace ledger with per-project violation and
complexity counts. The settings screen exposes the privacy controls
(offline mode, no cloud transmission, encrypted storage), the Docker
sandbox image and limits, the local model selection, and the analysis
depth and complexity threshold.

\subsection{Static Analysis}
The static-analysis screen (\cref{fig:static}) is the primary triage
surface. A file tree on the left is annotated with per-file finding
counts; the centre pane shows the selected source; the lower pane lists
that file's fused findings, each with a severity badge, the originating
modality, a code excerpt, and a suggested fix. \Cref{fig:static} shows
the model and pattern passes agreeing on a prompt-assembly weakness in
the case-study project.

\begin{figure}[t]
  \centering
  \includegraphics[width=\columnwidth]{static-analysis}
  \caption{Static analysis---file tree, source view, and per-file fused
           findings with remediation guidance.}
  \label{fig:static}
\end{figure}

\subsection{AI Review}
The AI-review screen (\cref{fig:aireview}) is an interactive companion to
the automated pass: the developer selects a file and asks the local model
free-form questions about it, or triggers a quick or deep review. All
exchanges stay on \texttt{127.0.0.1}.

\begin{figure}[t]
  \centering
  \includegraphics[width=\columnwidth]{ai-review}
  \caption{AI review---file-scoped conversation with the on-device
           \mbox{Llama~3.2} model.}
  \label{fig:aireview}
\end{figure}

\subsection{Vulnerabilities Register}
The vulnerabilities screen (\cref{fig:vulns}) is the cross-file view of
every fused finding: critical-count, model-corroborated-count and total
tiles above an expandable audit trail. Each row carries the file and line,
the category, a corroboration marker, and, on expansion, the code context
and remediation plan.

\begin{figure}[t]
  \centering
  \includegraphics[width=\columnwidth]{vulnerabilities}
  \caption{Vulnerabilities register---project-wide audit trail with
           severity and corroboration markers.}
  \label{fig:vulns}
\end{figure}

\subsection{Complexity}
The complexity screen (\cref{fig:complexity}) reports mean cyclomatic
complexity and decision-branch totals, a branching-trend chart, and a
ranked list of structural hot-spots with a short refactoring note for the
worst offender.

\begin{figure}[t]
  \centering
  \includegraphics[width=\columnwidth]{complexity}
  \caption{Complexity---mean cyclomatic complexity, branching trend, and
           ranked structural hot-spots.}
  \label{fig:complexity}
\end{figure}

\subsection{Dynamic Analysis and Container Insights}
The dynamic-analysis screen (\cref{fig:dynamic}) provisions the sandbox
and reports image build time, cold-start latency, live CPU, and memory
allocation, with the raw Docker engine log. The container-insights screen
overlays the static severity distribution and top-vulnerable-file ranking
with live telemetry when the sandbox is running. The build-and-CI screen
runs the project's build and test pipeline and streams its console.

\begin{figure}[t]
  \centering
  \includegraphics[width=\columnwidth]{dynamic-analysis}
  \caption{Dynamic analysis---sandbox build/start metrics, engine log,
           and memory allocation.}
  \label{fig:dynamic}
\end{figure}

\subsection{Risk Scoring and Reports}
The risk-scoring screen (\cref{fig:risk}) renders the \phindex\ as a gauge
with its band, the three sub-scores of \cref{sec:risk} as labelled bars,
and a short model-written rationale. The reports screen (\cref{fig:reports})
assembles the archival summary and finding ledger and exports the
paginated PDF or JSON described in \cref{sec:method}.

\begin{figure}[t]
  \centering
  \includegraphics[width=\columnwidth]{risk-scoring}
  \caption{Risk scoring---\phindex\ gauge and the three disclosed
           sub-scores of \cref{eq:sv,eq:sc,eq:sq}.}
  \label{fig:risk}
\end{figure}

\begin{figure}[t]
  \centering
  \includegraphics[width=\columnwidth]{reports}
  \caption{Reports---archival summary, reasoning ledger, and exportable
           finding register.}
  \label{fig:reports}
\end{figure}
